%@ pages 136-143
%@ status translated
% Original section 13, "Le champ associé à un préchamp", pp. 136--143.
%
% Translation notes (the mathematics is left exactly as printed):
%  * "n-précat (de Segal)" is rendered "Segal n-precategory", as in §2 and §3.
%  * The French prints "sous-$n+1$-cat'egorie de Segal" with a broken accent
%    command; read as "catégorie" and translated "Segal $n+1$-subcategory".
%  * The French prints "préfaiceaux" for "préfaisceaux"; translated
%    "presheaves".

\clearpage
\section{The stack associated to a prestack}
\label{associepage}

The construction of the sheaf associated to a presheaf plays a central role
in the theory of sheaves. The associated sheaf is defined by a universal
property. We begin by doing the same thing for the stack associated to a
prestack. In reading this section one must bear in mind that the analogue of
the sheaf/presheaf distinction is the distinction between the Segal
$n$-stacks for the topology $\Gg$ and the Segal $n$-stacks for the coarse
topology. The latter are simply the presheaves of Segal $n$-categories (the
only one of the conditions of \S9 that comes into play for the coarse
topology is the condition that each fibre be an $n$-category).

\begin{lemma}
\label{propuniv1}
Let $A$ be a Segal $n$-prestack and $A\rightarrow A'$ a choice of associated
stack. If $B$ is a $\Gg$-fibrant Segal $n$-prestack, then the morphism
$$
\underline{Hom}(A', B)\rightarrow \underline{Hom}(A,B)
$$
is an equivalence.
\end{lemma}
{\em Proof:} This follows from Proposition \ref{interne} for the internal
model category \newline
$nSePSh(\Xx ^{\Gg})$, since
$A\rightarrow A'$ is a $\Gg$-weak equivalence (by definition).
\eop

Recall that the objects of $nSeCHAMP (\Xx ^{\rm gro})$
are the Segal $n$-stacks
fibrant for the coarse topology, and that
$$
nSeCHAMP (\Xx ^{\Gg })\subset nSeCHAMP (\Xx ^{\rm gro})
$$
is the full enriched subcategory of the objects which are $\Gg$-fibrant,
which is
(by \ref{grofibchampimplfib}) also the full enriched subcategory
of the objects
which are stacks. In this context,
if $A\in nSeCHAMP (\Xx ^{\rm gro})_0$ is an object,
an {\em associated stack to $A$} is an
$\Gg$-fibrant replacement $A\rightarrow A'$ (one thus requires that $A'$ be
$\Gg$-fibrant, and not merely a stack, because one wants $A'$ to lie in
the subcategory above).
In other words, such an associated stack
is a morphism from $A$ to $A'$ in
$nSeCHAMP (\Xx ^{\rm gro})$, i.e. an object of the Segal
$n$-category $nSeCHAMP (\Xx ^{\rm gro})_{1/}(A,A')$.
In this framework we have the following universal property:

\begin{lemma}
\label{propuniv2}
Let $u:A\rightarrow A'$ be an associated stack in the sense of the preceding
paragraph.
For every $B\in nSeCHAMP (\Xx ^{\Gg})$, composition with
$u$ induces an equivalence
$$
nSeCHAMP (\Xx ^{\Gg})_{1/}(A',B)
\stackrel{\cong}{\rightarrow}
nSeCHAMP (\Xx ^{\rm gro})_{1/}(A,B).
$$
\end{lemma}
{\em Proof:}
In view of the definitions of the
$nPC$-enriched categories $nSeCHAMP(\ldots )$
and of the equality
$$
nSeCHAMP (\Xx ^{\Gg})_{1/}(A',B)=
nSeCHAMP (\Xx ^{\rm gro})_{1/}(A',B),
$$
this is a simple transcription of Lemma \ref{propuniv1}.
\eop



Consider now weak families of Segal $n$-categories over $\Xx ^o$,
which we regard as morphisms
$\Xx ^o\rightarrow nSeCAT'$ or as objects of
$\underline{Hom}(\Xx ^o, nSeCAT')$.
Recall that we have the equivalence
of Theorem \ref{correlation}
$$
nSeCHAMP(\Xx ^{\rm gro}) \cong \underline{Hom}(\Xx ^o, nSeCAT').
$$
If $F:\Xx ^o\rightarrow nSeCAT'$ we shall call {\em stack associated to
$F$} any pair $(F^{\rm ch}, f)$ where
$F^{\rm ch}:\Xx ^o\rightarrow nSeCAT'$ is a stack, i.e. satisfies the condition
of the second paragraph of \ref{correlation}, and where
$f: F \rightarrow F^{\rm ch}$ is, in the Segal $n+1$-category
$\underline{Hom}(\Xx ^o, nSeCAT')$, a
$1$-morphism
equivalent, via the equivalence of \ref{correlation} recalled
above, to a $\Gg$-weak equivalence $A\rightarrow A'$.
Under these conditions $F$ is equivalent to $A$ and $F^{\rm ch}$ is equivalent
to $A'$ (as objects of $\underline{Hom}(\Xx ^o, nSeCAT')$).
Moreover $A\rightarrow A'$ is an ``associated stack'' in the sense of the
preceding paragraph. This
allows us, starting from \ref{propuniv2}, to prove the
universal property of the following lemma.

Before stating it, recall that a $1$-morphism $f$ of
$\underline{Hom}(\Xx ^o, nSeCAT')$ is a morphism
$$
I=\Upsilon (\ast )\rightarrow \underline{Hom}(\Xx ^o, nSeCAT'),
$$
in other words, a morphism
$$
f:I\times \Xx ^o\rightarrow nSeCAT'.
$$
Its source is the restriction $f|_{\{ 0\} \times \Xx ^o}$ and its target
is the restriction $f|_{\{ 1\} \times \Xx ^o}$. If $f$ is a morphism with
source $F$ and target $F^{\rm ch}$ and if $G:\Xx ^o\rightarrow nSeCAT'$ is
another family, one can construct
an essentially well-defined morphism ``composition with $f$''
$$
-\circ f : \underline{Hom}(\Xx ^o, nSeCAT')_{1/}(F^{\rm ch}, G)
\rightarrow \underline{Hom}(\Xx ^o, nSeCAT')_{1/}(F, G).
$$
To give a precise definition of it, we denote by
$$
\underline{Hom}(\Xx ^o, nSeCAT')_{2/}(F,F^{\rm ch}, G;f)
$$
the fibre of
$$
r_{01}:\underline{Hom}(\Xx ^o, nSeCAT')_{2/}(F,F^{\rm ch}, G)
\rightarrow \underline{Hom}(\Xx ^o, nSeCAT')_{1/}(F,F^{\rm ch})
$$
over $f$; this fibre comes with two morphisms
(the first of which is an equivalence)
$$
r_{12}:\underline{Hom}(\Xx ^o, nSeCAT')_{2/}(F,F^{\rm ch}, G;f)
\stackrel{\cong}{\rightarrow}
\underline{Hom}(\Xx ^o, nSeCAT')_{1/}(F^{\rm ch}, G)
$$
and
$$
r_{02}:\underline{Hom}(\Xx ^o, nSeCAT')_{2/}(F,F^{\rm ch}, G;f)
\rightarrow
\underline{Hom}(\Xx ^o, nSeCAT')_{1/}(F, G).
$$
One obtains the morphism $-\circ f$ by choosing an inverse of $r_{12}$
up to equivalence (this is possible because all the Segal $n$-categories
considered are fibrant) and by composing it with $r_{02}$.

\begin{lemma}
\label{propuniv3}
Let $F\in \underline{Hom}(\Xx ^o, nSeCAT')_0$ and let $(F^{\rm ch}, f)$
be an associated stack in the sense defined above. Then for every stack
$G$, i.e. every object of $\underline{Hom}(\Xx ^o, nSeCAT')_0$ satisfying the
property of the second paragraph of \ref{correlation}, the morphism
``composition with $f$'' defined above is an equivalence of
Segal $n$-categories
$$
-\circ f : \underline{Hom}(\Xx ^o, nSeCAT')_{1/}(F^{\rm ch}, G)
\stackrel{\cong}{\rightarrow} \underline{Hom}(\Xx ^o, nSeCAT')_{1/}(F, G).
$$
\end{lemma}
{\em Proof:}
By definition, $f$ is equivalent to a morphism $u: A\rightarrow A'$
in the Segal $n+1$-category $nSeCHAMP(\Xx ^{\rm gro})$ which is an
associated stack for the topology $\Gg$, i.e. a $\Gg$-weak equivalence towards
a $\Gg$-fibrant object. One may also suppose, by \ref{correlation}, that $G$
comes from an object of $nSeCHAMP(\Xx ^{\Gg})$. The morphism ``composition with
$f$'' is then equivalent, via the equivalence of \ref{correlation}
$$
\underline{Hom}(\Xx ^o, nSeCAT')\cong
nSeCHAMP(\Xx ^{\rm gro}),
$$
to the morphism ``composition with $u$'' (the latter being well defined since
$nSeCHAMP(\Xx ^{\rm gro})$ is a category enriched over $nSeCat$). Lemma
\ref{propuniv2} allows us to conclude.
\eop

If $(F^{\rm ch,1}, f^1)$ and $(F^{\rm ch,2}, f^2)$ are two stacks associated
to $F$, let
$$
\underline{Hom}(\Xx ^o, nSeCAT')_{2/}(F,F^{\rm ch,1}, F^{\rm ch,2}; f^1, f^2)
:= (r_{01}, r_{02})^{-1}(f^1,f^2)
$$
be the fibre of
$$
(r_{01}, r_{02}): \underline{Hom}(\Xx ^o, nSeCAT')_{2/}(F,F^{\rm ch,1}, F^{\rm
ch,2}) \rightarrow
$$
$$
\underline{Hom}(\Xx ^o, nSeCAT')_{1/}(F,F^{\rm ch,1})
\times
\underline{Hom}(\Xx ^o, nSeCAT')_{1/}(F, F^{\rm ch,2})
$$
over $(f^1,f^2)$. One sees thanks to Lemma \ref{propuniv3} that
this fibre $(r_{01}, r_{02})^{-1}(f^1,f^2)$ is contractible, i.e. equivalent
(as a Segal $n$-category) to $\ast$. For this
one must observe that the
morphism $(r_{01}, r_{02})$ is a fibration of Segal $n$-precategories.
One thus obtains an essentially well-defined morphism between
$(F^{\rm ch,1}, f^1)$ and $(F^{\rm ch,2}, f^2)$.
One can show the same thing in the other direction
(and also the same thing for the two-dimensional diagrams
between $F^{\rm ch,1}$, $F^{\rm ch,2}$, $F^{\rm ch,1}$ or
$F^{\rm ch,2}$, $F^{\rm ch,1}$, $F^{\rm ch,2}$), which shows that the morphism
we have constructed
is an equivalence. This gives the essential uniqueness of the ``stack
associated to $F$''.

We now want to glue these associated stacks together into a ``weak''
functor
$$
nSeCHAMP(\Xx ^{\rm gro})\rightarrow nSeCHAMP(\Xx ^{\Gg}),
$$
in other words a functor to the fibrant replacement
$nSeCHAMP(\Xx ^{\Gg})'$. In fact, we shall construct a diagram
$$
\underline{Hom}(\Xx ^o, nSeCAT')\leftarrow C\rightarrow
\underline{Hom}(\Xx ^o, nSeCAT'),
$$
where the first morphism is an equivalence, and the second
takes its values in the
full Segal $n+1$-subcategory of stacks. This will give (by
\ref{correlation} and the usual arguments) a genuine morphism,
essentially well defined, to the desired fibrant replacement.
We shall use freely the notations $Diag$, cf.\ \S12, and $\Upsilon ^k$, cf.\
\cite{limits}.

We define $C$ as follows. The objects of $C$ are the objects of
\newline
$nSeCHAMP(\Xx ^{\rm gro})$, i.e. the Segal $n$-prestacks fibrant
for the coarse topology. For each object $F\in C_0$ we choose a
morphism
$\eta _F: F\rightarrow F'$ to a $\Gg$-fibrant replacement.
Then if $F_0,\ldots ,
F_p$ are objects, we define (with the notations of \S12)
$$
C_{p/}(F_0,\ldots , F_p):= Diag(\Xx ^o\times I^{(k)}\times I,
\Xx ^o\times \{ 0,\ldots , p\}
\times I;  \eta _{F_0}\sqcup \ldots \sqcup \eta _{F_p}; nSeCAT').
$$
This is the Segal $n$-category of diagrams in $nSeCAT'$ which have the shape
of a $p\times 1$ rectangle, restricting, for every $i$, to $\{ i\} \times I$
in the morphism $\eta _{F_i}: F_i\rightarrow F'_i$ regarded as a
morphism $I\rightarrow nSeCAT'$.

By dividing the rectangle, up to trivial cofibration, into a coproduct of
$1\times 1$ squares, and by using the fact that
$nSeCAT'$
is fibrant, one obtains that the Segal map
$$
C_{p/}(F_0,\ldots , F_p)\rightarrow
C_{1/}(F_0,F_1)\times \ldots \times C_{1/}(F_{p-1}, F_p)
$$
is an equivalence; hence $C$ is a Segal $n+1$-category.
Then, by dividing a $1\times 1$ square into two triangles,
one sees that a morphism
$$
E\rightarrow Diag(\Xx ^o\times I\times I,
\Xx ^o\times \{ 0,1\}\times I; \eta _{F_0}\sqcup \eta _{F_1};
nSeCAT')
$$
is nothing other than a pair of morphisms
$$
u: \Xx ^o\times \Upsilon ^2(\ast ,E)\rightarrow nSeCAT', \;\;\;
v: \Xx ^o\times \Upsilon  ^2(E,\ast )\rightarrow nSeCAT',
$$
with
$$
r_{01}(u)=\eta _{F_0},\;\; r_{12}(v)= \eta _{F_1}, \;\; r_{02}(u)=r_{02}(v).
$$
If we denote, for example, by $\underline{Hom}(\Xx ^o,nSeCAT')_{2/}(F_0, F'_0, F'_1;
\eta _{F_0},-)$ the fibre of
$$
r_{01}:\underline{Hom}(\Xx ^o,nSeCAT')_{2/}(F_0, F'_0, F'_1)\rightarrow
\underline{Hom}(\Xx ^o,nSeCAT')_{1/}(F_0, F'_0)
$$
over $\eta _{F_0}$, then we obtain
$$
Diag(\Xx ^o\times I\times I,
\Xx ^o\times \{ 0,1\}\times I; \eta _{F_0}\sqcup \eta _{F_1};
nSeCAT') =
$$
$$
\underline{Hom}(\Xx ^o,nSeCAT')_{2/}(F_0, F'_0, F'_1; \eta _{F_0},-)\times _{
\underline{Hom}(\Xx ^o,nSeCAT')_{1/}(F_0,  F'_1)}
$$
$$
\underline{Hom}(\Xx ^o,nSeCAT')_{2/}(F_0, F_1, F'_1; -,\eta _{F_1}).
$$

The universal property for the morphism $\eta _{F_0}$, and the fact that
$F'_1$ is $\Gg$-fibrant (i.e. a stack), imply that
the morphism
$$
u\mapsto r_{02}(u)
$$
is a trivial fibration of Segal $n$-precategories,
$$
r_{02}: \underline{Hom}(\Xx ^o,nSeCAT')_{2/}(F_0, F'_0, F'_1; \eta _{F_0},-)
\rightarrow \underline{Hom}(\Xx ^o,nSeCAT')_{1/}(F_0,  F'_1).
$$
On the other hand, the fact that
$$
\Upsilon (E)\sqcup ^{\{ 1\} }\Upsilon (\ast )
\rightarrow \Upsilon ^2(E,\ast )
$$
is a trivial cofibration implies that the morphism
$$
\underline{Hom}(\Xx ^o,nSeCAT')_{2/}(F_0, F_1, F'_1; -,\eta _{F_1})
\rightarrow
\underline{Hom}(\Xx ^o,nSeCAT')_{1/}(F_0, F_1)
$$
is an equivalence. One obtains that the morphism
$$
C_{1/}(F_0,F_1)=
Diag(\Xx ^o\times I\times I, \Xx ^o \times \{ 0,1\}\times I; \eta _{F_0}\sqcup
\eta _{F_1}; nSeCAT')
$$
$$
\rightarrow \underline{Hom}(\Xx ^o,nSeCAT')_{1/}(F_0,F_1)
$$
is an equivalence.

The restriction of the diagrams of $I^{(p)}\times I$ to
$I^{(p)}\times \{ 0\}$ provides a morphism
$$
C_{p/}(F_0,\ldots , F_p)\rightarrow\underline{Hom}(\Xx
^o,nSeCAT')_{p/}(F_0,\ldots , F_p),
$$
hence a morphism of Segal $n+1$-categories
$$
R_0: C\rightarrow \underline{Hom}(\Xx ^o,nSeCAT').
$$
We have shown that this morphism is fully faithful;
and it is essentially surjective by construction, thanks to
\ref{correlation}.

On the other hand, the restriction of the diagrams of $I^{(p)}\times I$ to
$I^{(p)}\times \{ 1\}$ provides a morphism
$$
C_{p/}(F_0,\ldots , F_p)\rightarrow\underline{Hom}(\Xx
^o,nSeCAT')_{p/}(F_0,\ldots , F_p),
$$
whence a morphism
$$
R_1: C\rightarrow \underline{Hom}(\Xx ^o,nSeCAT').
$$
These two morphisms form the desired diagram:
$$
\underline{Hom}(\Xx ^o,nSeCAT')
\stackrel{R_0}{\leftarrow}C
\stackrel{R_1}{\rightarrow}\underline{Hom}(\Xx ^o,nSeCAT').
$$


Observe that we have $R_0(F)=F$ and $R_1(F)=F'$.
We shall say that this diagram is the ``associated stack functor'',
also denoted ${\bf ch}$,
and we shall also write $F'= {\bf ch}(F)$.


It remains now to construct a natural transformation
between the identity and the functor constructed above. It
must be a diagram
$$
nSeCHAMP(\Xx ^{\rm gro})\stackrel{T_0}{\leftarrow} C\times I
\stackrel{T_1}{\rightarrow}
nSeCHAMP(\Xx ^{\rm gro}).
$$
We take $T_0:= R_0\circ pr_1$ and we want to construct a morphism $T_1$
which gives $R_0$ on $C\times \{ 0\}$ and $R_1$ on $C\times \{ 1\}$.
For this,
recall that the objects of $C\times I$ are the pairs $(F,\epsilon
)$ where $F$ is an object of $C$ and $\epsilon$ is $0$ or $1$. One
observes that
$$
(C\times I)_{p/}((F_0,\epsilon _0),\ldots , (F_p, \epsilon _p))
$$
is empty unless $\epsilon _0,\ldots , \epsilon _p$ is a
non-decreasing sequence, in which case it is simply
$C_{p/}(F_0,\ldots , F_p)$. To every non-decreasing sequence $\epsilon
_0,\ldots , \epsilon _p$ one associates the morphism
$$
i_{\epsilon} :I^{(p)}\rightarrow I^{(p)}\times I
$$
defined by $ i_{\epsilon}(k):= (k, \epsilon _k)$.

The inverse image under this morphism
therefore provides, for each sequence $\epsilon$, a morphism
$$
i_{\epsilon}^{\ast}: C_{p/}(F_0,\ldots , F_p)\rightarrow \underline{Hom}(\Xx
^o,nSeCAT')_{p/}(F_0^{(\epsilon _0)},\ldots , F_p^{(\epsilon _p)} ).
$$
Here $F^{(0)}$ and $F^{(1)}$ denote
$F$ and $F'$ respectively.

For $\epsilon = (0,\ldots , 0)$,
one recovers the restriction morphism that we
used to define $R_0$, while
for $\epsilon =(1,\ldots , 1)$ one recovers the one that we
used to define $R_1$. For every sequence
of objects $(F_0,\epsilon _0),\ldots , (F_p, \epsilon _p)$ one
defines, with the help
of $i_{\epsilon}^{\ast}$, the morphism
$$
(C\times I)_{p/}((F_0,\epsilon _0),\ldots , (F_p, \epsilon _p))\rightarrow
\underline{Hom}(\Xx
^o,nSeCAT')_{p/}(F_0^{(\epsilon _0)},\ldots , F_p^{(\epsilon _p)} );
$$
which allows one to define the desired morphism
$$
T_1: C\times I \rightarrow \underline{Hom}(\Xx
^o,nSeCAT').
$$

One can check that $T_1|_{\{ F\} \times I}$ is indeed
the morphism $F\rightarrow F'$.
The diagram $(T_0,T_1)$ therefore indeed constitutes a natural transformation
$$
1\rightarrow {\bf ch}
$$
between a prestack and its associated stack.

\begin{theorem}
\label{cestadjoint}
With this natural transformation, the ``associated stack'' functor
${\bf ch}$ is homotopically adjoint to the inclusion
$$
nSeCHAMP(\Xx ^{\Gg})\subset nSeCHAMP(\Xx ^{\rm gro}).
$$
\end{theorem}
{\em Proof:}
The necessary equivalence is a consequence of the universal property
of ${\bf ch}(F)$, \ref{propuniv2} for example.
\eop

\subnumero{Computing the morphisms in the associated stack}

We have the following formula:

\begin{lemma}
\label{formule}
For a Segal $n$-prestack $F$ whose
values are Segal $n$-categories ($n\geq 1$), and for $x,y\in F(X)$,
there is an object-by-object equivalence over $\Xx /X$,
$$
({\bf ch}(F))_{1/}(x,y) \cong {\bf ch}(F_{1/}(x,y)).
$$
\end{lemma}
{\em Proof:}
Note that
$({\bf ch}(F))_{1/}(x,y)$ is a Segal $n-1$-stack by \ref{critere}.
Since $F\rightarrow {\bf ch}(F)$ is a $\Gg$-weak equivalence, by
definition (\S3)
$$
F_{1/}(x,y)\rightarrow ({\bf ch}(F))_{1/}(x,y)
$$
is a $\Gg$-weak equivalence. Then by definition (\S9)
$({\bf ch}(F))_{1/}(x,y)$ is the Segal $n-1$-stack associated to
$F_{1/}(x,y)$.
\eop



\subnumero{Elementary version}
It seems difficult to obtain the functor ${\bf ch}$ directly by using the
structure of internal model category, because the ``$\Gg$-fibrant
replacement'' functor is not strictly compatible with direct products.
One can however remark that the Dwyer--Kan localisation theory
immediately gives
$$
L(nSePCh , W^{\rm gro}) \rightarrow L(nSePCh, W^{\Gg}),
$$
where $W^{\rm gro}\subset nSePCh$ is the subcategory
of weak equivalences for the coarse topology, and
$W^{\Gg}\subset W^{\rm gro}$ is its subcategory of $\Gg$-weak
equivalences.
The fact that $W^{\rm gro}$ is contained in $W^{\Gg}$ gives directly the
morphism above.
Now the result of Theorem
\ref{intereqloc}
shows that the preceding morphism is equivalent to a morphism
$$
nSeCHAMP(\Xx ^{\rm gro})^{int, 1}\rightarrow
nSeCHAMP'(\Xx ^{\Gg})^{int, 1}
$$
(we have taken a fibrant replacement denoted $nSeCHAMP'$ at the target).
This morphism is the restriction of ${\bf ch}$ to the $1$-groupic interior.
In the case $n=0$, i.e. of simplicial presheaves, one has
$$
0SeCHAMP(\Xx ^{\rm gro})^{int, 1}=0SeCHAMP(\Xx ^{\rm gro})
$$
and
$$
0SeCHAMP(\Xx ^{\Gg})^{int, 1}=0SeCHAMP(\Xx ^{\Gg})
$$
and one thus obtains (i.e. with the Dwyer--Kan localised category) an elementary
definition of the functor ${\bf ch}$.
